\documentclass[11pt]{article}
\usepackage[a4paper,margin=2.5cm]{geometry}
\usepackage{amsmath,amssymb}
\usepackage{graphicx}
\usepackage{booktabs}
\usepackage{longtable}
\usepackage{array}
\usepackage{hyperref}
\usepackage{enumitem}
\usepackage{xcolor}
\usepackage{caption}
\usepackage{microtype}
\sloppy
\graphicspath{{figures/}}
\setlist[itemize]{leftmargin=1.2em}
\setlist[enumerate]{leftmargin=1.4em}
\renewcommand{\arraystretch}{1.15}

\title{BOAMP M1 Buyer-Identity Enrichment\\
\large Consolidated Technical Report for SIREN Enrichment, Historical Alias
Propagation, Profile-Domain Audit, and Survival Sensitivity}
\author{Gigalis Predictive-Modeling Internship --- Current Independent Study}
\date{Repository state of July 2026; M1 outputs generated 2026-07-15}

\begin{document}
\maketitle

\begin{abstract}
This report is the consolidated technical reference for the M1
buyer-identity enrichment branch. M1 is not the official baseline. The
official study result remains M0, the BOAMP-only balanced specification.
M1 is a sensitivity analysis designed to answer one controlled question:
\emph{how much does buyer-name fragmentation in M0 affect candidate
coverage and proxy-recurrence counts when buyer identity is enriched with
validated SIREN evidence?}

M1 keeps the M0 eligible source population, temporal window logic,
TF--IDF object-text score, CPV score, temporal score, ranking rule, and
balanced threshold. The controlled change is buyer reconciliation. M1
uses exact BOAMP SIRET when present, then a validated SIREN from either a
direct notice-level enrichment join or a conservative historical
name--department alias bridge, and finally the original M0 normalized-name
fallback. The M0 balanced threshold is reused unchanged
($0.323022$), so M1's higher event count is not produced by relaxing the
decision threshold.

On the same 3,159 eligible digital/ICT source notices, M1 increases
blocking coverage from 39.1\% to 47.6\%, candidate pairs from 6,137 to
8,228, and accepted balanced proxy-recurrence links from 618 to 847. The
overall proxy-recurrence rate rises from 19.6\% to 26.8\%. This is useful
evidence that buyer identifiers matter, but it is not proof that the
incremental links are legal renewals. M1 adds 273 links not present in M0,
loses 44 M0 links under the new ranking, and has link-set Jaccard overlap
0.644 with M0. Incremental links have weaker text evidence than unchanged
M0 links (median TF--IDF similarity 0.078 versus 0.175), so M1 should
remain a sensitivity specification until manual validation is completed.

A procurement-profile domain audit is included as a second-stage evidence
overlay. It does not create buyer identifiers and does not add links. It
classifies profile domains as supporting, neutral, weak, ambiguous, or
missing evidence for the M1 alias decisions. The profile-audited variant
removes zero links and is byte-for-byte identical in headline counts to
M1: 847 accepted links, 47.6\% blocking coverage, 26.8\% linking rate, and
Jaccard 1.000 versus M1.
\end{abstract}

\tableofcontents

\section{Purpose and Position}
\label{sec:purpose}

The M1 branch exists because M0 deliberately avoids external enrichment.
That makes M0 clean and auditable, but it also means that buyer identity
depends only on whatever BOAMP provides in the notice: SIRET if present,
SIREN if present, otherwise a normalized buyer name. In practice, many
digital/ICT notices do not carry a usable identifier, and the same public
buyer can appear under multiple name variants.

M1 is therefore a sensitivity analysis, not a replacement. Its role is:

\begin{enumerate}
  \item preserve the official M0 design as the reference result;
  \item change only buyer reconciliation and candidate blocking;
  \item quantify how many candidate pairs and proxy-recurrence links are
        recovered when missing buyer identifiers are partially repaired;
  \item make every enriched identity and alias decision traceable; and
  \item separate useful sensitivity evidence from validated ground truth.
\end{enumerate}

The current repository decision is:

\begin{quote}
M0 balanced is the main result and official survival handoff. M1
buyer-SIREN enrichment is a sensitivity analysis. The event remains a
constructed proxy recurrence, not a verified legal renewal.
\end{quote}

\section{Input Data and Lineage}
\label{sec:lineage}

\subsection{M0 reference inputs}

M1 reads the M0 outputs as fixed reference inputs. The source population
and survival eligibility are not recomputed from scratch under a different
scope. The principal M0 inputs are:

\begin{itemize}
  \item \texttt{data/processed/boamp\_m0\_sources.csv};
  \item \texttt{data/processed/boamp\_m0\_candidate\_pairs.csv};
  \item \texttt{data/processed/boamp\_m0\_links\_balanced.csv};
  \item \texttt{data/processed/boamp\_survival\_m0\_balanced.csv}.
\end{itemize}

The M1 consistency audit records SHA-256 traces for these M0 files and
checks that they remain read-only reference inputs. This is important:
M1 does not silently rewrite M0.

\subsection{SIREN enrichment dataset}

The direct buyer enrichment source is the Hugging Face dataset
\texttt{Data-Gouv-ML/jointure-boamp-siren-cote-acheteurs-2024-2025-et-2026},
commit \texttt{4bff9b1c5d2f4ad12834174e042828b7e52013d9}, downloaded on
2026-07-15. The join key is:

\[
\texttt{B\_17\_idweb} \longrightarrow \texttt{notice\_id}.
\]

The live audit found no duplicate notice identifiers in the three Parquet
files, so the enrichment merge is one-to-one on notice ID for the years it
covers.

\begin{table}[h]
\centering
\caption{Notice-level SIREN enrichment files used by M1.}
\begin{tabular}{lrrrr}
\toprule
Year file & Rows & Date min & Date max & Overlap with M0 sources\\
\midrule
2024 & 136,625 & 2024-01-01 & 2024-12-31 & 210\\
2025 & 137,423 & 2025-01-01 & 2025-12-31 & 189\\
2026 & 49,890 & 2026-01-01 & 2026-05-28 & 59\\
\midrule
Total & 323,938 & -- & -- & 458\\
\bottomrule
\end{tabular}
\end{table}

Across the cleaned BOAMP table, 18,098 notices overlap the enrichment
dataset. Among rows where BOAMP provides a SIRET-derived SIREN and the
external enrichment also provides a valid SIREN, the agreement rate is
87.4\% (6,937 agreements out of 7,939 rows). The 1,002 disagreements are
not discarded silently; they are exported to
\texttt{reports/tables/buyer\_identifier\_conflicts.csv} and flagged in
the M1 identity provenance columns.

\subsection{Procurement-profile dataset}

The procurement-profile audit uses the data.gouv.fr dataset ``Liste des
profils d'acheteurs par entites adjudicatrices'', downloaded as
\texttt{data/raw/procurement\_profiles/profils-acheteurs.csv} on
2026-07-15. The producer is Donn\'eesPubliques.org and the recorded
license is \texttt{lov2}.

This dataset is evidence-only. Its row unit is a buyer-name /
profile-domain / participation-domain summary row with a latest observed
BOAMP date. It has no BOAMP notice ID, no SIREN, no SIRET, no city, and no
postcode. Therefore M1 never uses procurement-profile domains as legal
buyer identifiers and never creates a candidate pair from profile evidence.

\begin{table}[h]
\centering
\caption{Procurement-profile dataset audit.}
\begin{tabular}{lr}
\toprule
Metric & Value\\
\midrule
Rows & 23,260\\
Columns & 5\\
Observed date range & 2015-03-02 to 2020-12-31\\
Distinct normalized buyer names & 15,039\\
Distinct normalized profile domains & 4,700\\
Distinct participation domains & 900\\
Malformed profile URL rows & 396\\
Malformed participation URL rows & 22\\
Exact duplicate buyer/profile/participation rows & 1,715\\
\bottomrule
\end{tabular}
\end{table}

\section{End-to-End M1 Data Preprocessing}
\label{sec:m1-end-to-end-preprocessing}

This section records the full preprocessing path from fixed M0 inputs to
the M1-ready source table. It is intentionally written as a data lineage,
not as an argument that M1 is superior to M0. The official M0 files remain
fixed reference inputs throughout.

\subsection{Step 0: freeze the M0 reference surface}

M1 starts after the M0 preprocessing, candidate generation, linkage, and
survival handoff already exist. The M1 script reads the M0 source table,
candidate-pair table, balanced link table, and balanced survival handoff to
recover the exact source population, score threshold, and comparison
baseline. It does not recompute the BOAMP download scope, the digital/ICT
filter, M0 duration imputation, or the official M0 balanced links. This
design means every M1 difference can be interpreted as a buyer-identity
sensitivity rather than a hidden change in the study population.

\subsection{Step 1: load and audit the external SIREN files}

\texttt{scripts/build\_m1\_buyer\_siren\_experiment.py} first loads the
three Parquet files under
\texttt{data/raw/buyer\_siren\_enrichment\_m1/}. Only the documented
BOAMP and Sirene columns listed in \texttt{ENRICHMENT\_COLUMNS} are read.
For each file the script records row counts, date coverage, missing notice
IDs, duplicate notice IDs, missing SIREN/SIRET rates, file size, and SHA-256
hash. The audit is written to
\texttt{reports/tables/buyer\_siren\_enrichment\_audit.csv} and
\texttt{buyer\_siren\_enrichment\_audit.json}. The live audit supports the
join key \texttt{B\_17\_idweb -> notice\_id}: no duplicated notice IDs were
found in the 2024, 2025, or 2026 enrichment files.

\subsection{Step 2: merge notice-level enrichment without overwriting M0}

The enrichment rows are renamed to explicit M1 fields and left-joined to
\texttt{data/processed/boamp\_clean\_m0\_no\_enrichment.csv} by
\texttt{notice\_id}. The result is an all-notice enriched table, not yet the
source-only analysis table. Dates are re-parsed after the merge so that
subsequent temporal logic does not inherit object/string date types from the
join. BOAMP buyer identifiers remain in their original fields; external
SIREN/SIRET values are stored separately as enriched fields. This separation
is central to the audit trail: M1 can use external evidence, but it does not
rewrite BOAMP-provided SIRET/SIREN values.

\subsection{Step 3: clean identifiers and classify provenance}

The preprocessing then normalizes buyer names and cleans identifier fields.
SIREN values must be nine digits; SIRET values must be fourteen digits; and
validated external SIREN values must pass the identifier checksum before
direct use. The script derives:

\begin{itemize}
  \item BOAMP effective SIREN from BOAMP SIREN or from the first nine digits
        of a valid BOAMP SIRET;
  \item valid enriched SIREN and recovered enriched SIRET from the external
        notice-level join;
  \item buyer-name and department keys used only for alias auditing;
  \item disagreement flags where BOAMP SIRET-derived SIREN and enriched
        SIREN are both present but disagree.
\end{itemize}

Identifier conflicts are not resolved silently. They are exported to
\texttt{reports/tables/buyer\_identifier\_conflicts.csv} and reflected in
the identity-source and confidence columns.

\subsection{Step 4: build the historical alias bridge}

Direct enrichment only covers recent BOAMP notices. To test whether older
name-only rows are fragmented, M1 builds a conservative historical bridge
from directly enriched, non-conflicting rows. The bridge groups by normalized
buyer name and department key. A historical alias is eligible for automatic
propagation only when one valid enriched SIREN is uniquely observed for the
name--department cell, the cell has no BOAMP conflict, the name is not a
generic institutional label, and at least two notices support the mapping.
Ambiguous cells are retained in
\texttt{reports/tables/buyer\_alias\_ambiguous.csv} but are not propagated.
The accepted bridge is written to
\texttt{reports/tables/buyer\_alias\_bridge.csv}.

\subsection{Step 5: construct M1 buyer identity fields}

For every cleaned BOAMP notice, M1 applies the identity priority rule:
preserve a BOAMP SIRET-derived identity first; otherwise use a valid direct
enriched SIREN when there is no BOAMP conflict; otherwise use an automatic
historical alias only for old rows with no BOAMP identifier; otherwise fall
back to the original normalized buyer-name key. The output fields record the
resulting M1 buyer key, the identity source, confidence label, direct versus
inferred status, and conflict status. The source-only M1 table is then formed
by selecting the same 3,159 source notices used by M0 and attaching those M1
identity/provenance fields. This table is written to
\texttt{data/processed/boamp\_clean\_m1\_buyer\_enriched.csv}; the all-notice
enriched table is written to
\texttt{data/processed/boamp\_clean\_m1\_buyer\_enriched\_all\_notices.csv}.

\subsection{Step 6: regenerate candidate pairs under the M1 buyer keys}

Only after identity preprocessing is complete does M1 generate candidate
pairs. The candidate universe keeps M0's source population, estimated end
dates, temporal window, TF--IDF text representation, CPV hierarchy score,
temporal score, composite weights, maximum of 30 candidates per source, and
balanced threshold. The controlled change is buyer compatibility. A source
and later candidate can enter the M1 candidate table through exact SIRET,
same SIREN, historical alias, or original M0 name fallback. The resulting
candidate table is
\texttt{data/processed/boamp\_m1\_candidate\_pairs.csv}.

\subsection{Step 7: create links, survival rows, and diagnostics}

M1 ranks candidates within each source by the M1 composite score, accepts
the rank-1 candidate when its score is at least the M0 balanced threshold,
and writes \texttt{data/processed/boamp\_m1\_links\_balanced.csv}. The
survival handoff keeps one row per eligible source and writes
\texttt{data/processed/boamp\_survival\_m1\_balanced.csv}. Downstream audit
tables compare M0 and M1 at source and link level, identify incremental M1
links, summarize recovery mechanisms, measure candidate reuse, produce
negative-control diagnostics, and export an unlabeled manual-review sample.

\begin{table}[h]
\centering
\caption{End-to-end M1 preprocessing products.}
\scriptsize
\begin{tabular}{p{0.40\linewidth}p{0.48\linewidth}}
\toprule
Output & Role in the M1 preprocessing lineage\\
\midrule
\path{buyer_siren_enrichment_audit.csv} & File-level audit of external SIREN/SIRET enrichment inputs.\\
\path{buyer_identifier_conflicts.csv} & Rows where BOAMP-derived and enriched legal identifiers disagree.\\
\path{buyer_alias_bridge.csv} & Automatically eligible historical name--department aliases.\\
\path{buyer_alias_ambiguous.csv} & Ambiguous or review-only aliases not propagated automatically.\\
\path{boamp_clean_m1_buyer_enriched_all_notices.csv} & All cleaned BOAMP notices after notice-level enrichment and provenance assignment.\\
\path{boamp_clean_m1_buyer_enriched.csv} & The same 3,159 eligible source notices as M0, with M1 identity fields attached.\\
\path{boamp_m1_candidate_pairs.csv} & M1 candidate pairs generated after buyer-identity preprocessing.\\
\path{boamp_m1_links_balanced.csv} & Accepted M1 proxy-recurrence links at the M0 balanced threshold.\\
\path{boamp_survival_m1_balanced.csv} & One-row-per-source M1 survival handoff with provenance columns.\\
\bottomrule
\end{tabular}
\end{table}

\subsection{Step 8: profile-domain audit as a separate overlay}

The procurement-profile step is not part of the core SIREN preprocessing.
It reads M1 aliases and M1 candidate pairs after the buyer-enrichment branch
has already been built. It normalizes profile domains, classifies them as
buyer-specific, highly concentrated, shared institutional, generic platform,
ambiguous, or malformed, and annotates M1 alias/candidate evidence. It may
flag conflicts for manual review, but it does not create legal identifiers,
does not create new buyer matches, and does not replace M0 or M1.

\section{Notation}
\label{sec:notation}

For each eligible source or candidate notice \(i\), define:

\begin{itemize}
  \item \(d_i\): publication date;
  \item \(\hat{e}_i\): estimated end date inherited from M0;
  \item \(u_i\): normalized buyer name;
  \item \(t_i\): BOAMP SIRET, if format-valid and present;
  \item \(\sigma_i^{B}\): BOAMP effective SIREN, using BOAMP SIREN if
        present or the first nine digits of a valid BOAMP SIRET;
  \item \(\sigma_i^{E}\): valid enriched SIREN from the direct
        notice-level enrichment join;
  \item \(\sigma_i^{A}\): SIREN inferred from an automatic
        name--department historical alias;
  \item \(k_i^{0}\): original M0 buyer key;
  \item \(k_i^{1}\): M1 primary buyer key;
  \item \(y_i\): publication year.
\end{itemize}

Let \(H_i\) denote the historical-alias eligibility condition:

\[
H_i =
\left[
y_i \le 2023
\right]
\land
\left[
\sigma_i^{A} \ne \varnothing
\right]
\land
\left[
t_i = \varnothing
\right]
\land
\left[
\sigma_i^{B} = \varnothing
\right].
\]

So a historical alias can only fill a missing identifier for old notices.
It does not overwrite a BOAMP SIRET or BOAMP SIREN.

\section{M1 Buyer Identity Construction}
\label{sec:identity}

\subsection{Direct enrichment fields}

The direct enrichment process first cleans the external SIREN/SIRET fields:
non-digits are removed, SIREN must be nine digits, SIRET must be fourteen
digits, and external SIREN must pass the checksum validation before it is
considered valid for direct use.

The implementation preserves three separate ideas:

\begin{itemize}
  \item BOAMP identifiers, which are the original source of truth for M0;
  \item enriched identifiers, which are external evidence joined by
        notice ID for 2024--2026 notices;
  \item conflict flags, which record disagreement between BOAMP SIRET-derived
        SIREN and enriched SIREN.
\end{itemize}

This separation is what prevents M1 from becoming an untraceable overwrite
of M0.

\subsection{Historical alias bridge}

The historical alias bridge is built from directly enriched, non-conflicting
notices. For each normalized buyer name and department key, M1 examines the
set of observed valid enriched SIRENs. An alias is automatically propagated
only if:

\begin{enumerate}
  \item exactly one SIREN is observed for the name--department cell;
  \item the cell has no BOAMP-SIRET conflict;
  \item the buyer name is not a generic label such as \texttt{mairie},
        \texttt{commune}, or \texttt{centre hospitalier};
  \item the support count is at least two notices.
\end{enumerate}

The bridge contains 1,247 automatically eligible aliases and 1,912
ambiguous or review-only aliases. Ambiguous aliases are retained for audit
but are not propagated.

\subsection{Piecewise identity rule}

The M1 SIREN value is constructed as:

\[
\sigma_i^{1} =
\begin{cases}
\sigma_i^{B}, & \text{if BOAMP effective SIREN exists},\\
\sigma_i^{E}, & \text{else if valid enriched SIREN exists and no BOAMP conflict},\\
\sigma_i^{A}, & \text{else if } H_i \text{ is true},\\
\varnothing, & \text{otherwise.}
\end{cases}
\]

The M1 primary buyer key is then:

\[
k_i^{1} =
\begin{cases}
\texttt{SIRET:}t_i, & \text{if BOAMP SIRET } t_i \text{ exists},\\
\texttt{SIREN:}\sigma_i^{1}, & \text{else if } \sigma_i^{1} \text{ exists},\\
\texttt{NAME:}u_i, & \text{else if normalized buyer name } u_i \text{ exists},\\
\varnothing, & \text{otherwise.}
\end{cases}
\]

This rule is implemented in
\texttt{scripts/build\_m1\_buyer\_siren\_experiment.py} in
\texttt{apply\_m1\_identity}. The exact practical consequence is that M1
can move some M0 \texttt{NAME:...} rows into \texttt{SIREN:...} keys, but
it does not replace an available BOAMP SIRET with an external value.

\begin{table}[h]
\centering
\caption{M1 identity-source distribution among eligible sources.}
\begin{tabular}{lr}
\toprule
Identity source & Source notices\\
\midrule
\texttt{NAME\_FALLBACK} & 1,231\\
\texttt{UNIQUE\_NAME\_DEPARTMENT\_ALIAS} & 919\\
\texttt{DIRECT\_BOAMP\_SIRET} & 686\\
\texttt{DIRECT\_ENRICHMENT\_NOTICE\_JOIN} & 296\\
\texttt{DIRECT\_BOAMP\_SIRET\_CONFLICT\_FLAGGED} & 27\\
\bottomrule
\end{tabular}
\end{table}

\subsection{Confidence labels}

The identity confidence labels are deliberately descriptive rather than
probabilistic:

\begin{itemize}
  \item \texttt{HIGH}: direct BOAMP SIRET, including conflict-flagged SIRET
        rows where BOAMP is still preserved;
  \item \texttt{MEDIUM\_HIGH}: direct valid enriched SIREN joined by notice ID;
  \item \texttt{MEDIUM}: direct BOAMP SIREN without SIRET;
  \item \texttt{MEDIUM\_LOW}: unique name--department historical alias;
  \item \texttt{LOW}: normalized-name fallback;
  \item \texttt{CONFLICT}: unresolved enriched identifier disagreement.
\end{itemize}

In the current eligible source population, the counts are 713 high,
296 medium-high, 919 medium-low, and 1,231 low.

\section{Candidate Generation and Scoring}
\label{sec:candidates}

\subsection{What stays identical to M0}

M1 keeps the M0 design choices below:

\begin{itemize}
  \item same eligible source population: 3,159 digital/ICT
        \texttt{APPEL\_OFFRE} source notices;
  \item same estimated end date and duration-imputation logic;
  \item same TF--IDF vectorization for object text;
  \item same CPV hierarchy score;
  \item same temporal score centered on the source notice's estimated end
        date;
  \item same maximum of 30 candidates per source after temporal proximity
        ordering;
  \item same composite-score weights;
  \item same M0 balanced threshold: \(0.323022\).
\end{itemize}

The current M1 temporal window is six months. It is computed from the same
median-duration rule used by the implementation and clipped to the
configured floor/cap interval \([6,24]\) months.

\subsection{Buyer compatibility mechanisms}

For a source notice \(i\) and later candidate notice \(j\), M1 materializes
a candidate only if at least one buyer-compatibility mechanism is satisfied:

\[
m_{ij} =
\begin{cases}
\texttt{EXACT\_SIRET}, & t_i=t_j \ne \varnothing,\\
\texttt{HISTORICAL\_ALIAS}, & \sigma_i^{1}=\sigma_j^{1}\ne\varnothing
  \text{ and either side uses an alias},\\
\texttt{SAME\_SIREN}, & \sigma_i^{1}=\sigma_j^{1}\ne\varnothing,\\
\texttt{M0\_NAME}, & k_i^{0}=k_j^{0}\ne\varnothing,\\
\varnothing, & \text{otherwise.}
\end{cases}
\]

If the same candidate is reachable through several mechanisms, M1 keeps
the highest-priority mechanism in this order:

\[
\texttt{EXACT\_SIRET}
>
\texttt{SAME\_SIREN}
>
\texttt{HISTORICAL\_ALIAS}
>
\texttt{M0\_NAME}.
\]

The buyer-mechanism score is:

\begin{table}[h]
\centering
\caption{Buyer mechanism scores used in the M1 composite score.}
\begin{tabular}{lr}
\toprule
Mechanism & \(s_{\mathrm{buyer}}\)\\
\midrule
\texttt{EXACT\_SIRET\_RECONCILIATION} & 1.00\\
\texttt{SAME\_SIREN\_RECONCILIATION} & 0.90\\
\texttt{HISTORICAL\_ALIAS\_RECONCILIATION} & 0.75\\
\texttt{M0\_NAME\_FALLBACK} & 0.60\\
\bottomrule
\end{tabular}
\end{table}

\subsection{Composite score}

For a candidate pair \((i,j)\), M1 uses:

\[
S_{ij}^{1}
=
0.35\,s_{\mathrm{text},ij}
+0.30\,s_{\mathrm{cpv},ij}
+0.25\,s_{\mathrm{time},ij}
+0.10\,s_{\mathrm{buyer},ij}.
\]

The temporal component is:

\[
s_{\mathrm{time},ij}
=
\max\left(1-\frac{|d_j-\hat{e}_i|}{w},\,0\right),
\]

where \(w=6\) months in the current run. Candidate \(j\) must be later
than source \(i\), and it must fall within the temporal window around the
source's estimated end date. M1 then ranks candidates within each source by
\(S_{ij}^{1}\), keeps the top-ranked candidate, and accepts it if:

\[
S_{ij}^{1} \ge 0.323022.
\]

\section{M1 Outputs}
\label{sec:outputs}

The principal M1 output files are:

\begin{itemize}
  \item \texttt{data/processed/boamp\_clean\_m1\_buyer\_enriched.csv};
  \item \texttt{data/processed/boamp\_m1\_candidate\_pairs.csv};
  \item \texttt{data/processed/boamp\_m1\_links\_balanced.csv};
  \item \texttt{data/processed/boamp\_survival\_m1\_balanced.csv};
  \item \texttt{reports/tables/m0\_m1\_linkage\_comparison.csv};
  \item \texttt{reports/tables/m0\_m1\_link\_status\_changes.csv};
  \item \texttt{reports/tables/m1\_incremental\_links.csv};
  \item \texttt{reports/tables/m1\_manual\_validation\_sample\_unlabeled.csv}.
\end{itemize}

The survival handoff remains one row per eligible source. Compared with
M0, it adds M1 provenance columns such as
\texttt{buyer\_key\_m1\_primary}, \texttt{buyer\_identity\_source},
\texttt{buyer\_identity\_confidence}, and
\texttt{buyer\_match\_mechanism}.

\section{M0 versus M1 Results}
\label{sec:results}

\begin{table}[h]
\centering
\caption{Headline comparison between official M0 and M1 sensitivity.}
\begin{tabular}{lrr}
\toprule
Metric & M0 official & M1 sensitivity\\
\midrule
Eligible source notices & 3,159 & 3,159\\
Sources with at least one candidate & 1,236 & 1,504\\
Zero-candidate sources & 1,923 & 1,655\\
Blocking coverage & 39.1\% & 47.6\%\\
Candidate pairs & 6,137 & 8,228\\
Accepted balanced links & 618 & 847\\
Overall proxy-recurrence rate & 19.6\% & 26.8\%\\
Conditional acceptance among sources with candidates & 50.0\% & 56.3\%\\
Unique selected candidates & 443 & 570\\
Candidate reuse rate & 27.3\% & 29.5\%\\
Maximum candidate multiplicity & 6 & 8\\
\bottomrule
\end{tabular}
\end{table}

The M1 increase is mostly a blocking effect. M1 recovers 269 sources that
had no M0 candidate but receive at least one M1 candidate. At the final
accepted-link level:

\begin{itemize}
  \item 574 links are identical in M0 and M1;
  \item 273 accepted M1 links are not present in M0;
  \item 44 M0 links are not retained under M1 ranking/thresholding;
  \item 43 sources are linked in both methods but to a different selected
        candidate;
  \item link-set Jaccard overlap is 0.644.
\end{itemize}

\subsection{Accepted links by mechanism}

\begin{table}[h]
\centering
\caption{M1 accepted links by buyer-match mechanism.}
\begin{tabular}{lrr}
\toprule
Mechanism & Candidate pairs & Accepted links\\
\midrule
\texttt{HISTORICAL\_ALIAS\_RECONCILIATION} & 3,537 & 383\\
\texttt{M0\_NAME\_FALLBACK} & 3,106 & 254\\
\texttt{EXACT\_SIRET\_RECONCILIATION} & 883 & 120\\
\texttt{SAME\_SIREN\_RECONCILIATION} & 702 & 90\\
\bottomrule
\end{tabular}
\end{table}

The historical-alias mechanism is the largest accepted-link source. This
is exactly why M1 cannot be promoted to the official baseline without
manual validation: the main incremental mechanism is useful, but weaker
than direct BOAMP identifiers.

\section{Incremental-Link Diagnostics}
\label{sec:incremental}

M1's incremental links should be read carefully. They are not random new
events; they are links made possible by better buyer reconciliation. But
their internal diagnostics are weaker than the unchanged M0 links.

\begin{table}[h]
\centering
\caption{Incremental M1 links by recovery mechanism.}
\small
\begin{tabular}{lrrrrr}
\toprule
Recovery mechanism & Links & Median score & Median margin & Median text & CPV mismatch\\
\midrule
Historical alias & 198 & 0.360 & 0.066 & 0.073 & 25.8\%\\
Same SIREN & 75 & 0.388 & 0.073 & 0.085 & 14.7\%\\
\bottomrule
\end{tabular}
\normalsize
\end{table}

\begin{table}[h]
\centering
\caption{Score diagnostics: unchanged M0 links versus incremental M1 links.}
\begin{tabular}{lrr}
\toprule
Median diagnostic & Unchanged M0 links in M1 & Incremental M1 links\\
\midrule
Composite score & 0.407 & 0.372\\
TF--IDF text similarity & 0.175 & 0.078\\
CPV score & 0.100 & 0.100\\
Temporal score & 0.825 & 0.836\\
Buyer mechanism score & 0.750 & 0.750\\
Top1--top2 margin & 0.080 & 0.070\\
\bottomrule
\end{tabular}
\end{table}

The key interpretation is:

\[
\text{M1 improves candidate visibility}
\;\not\Rightarrow\;
\text{M1 proves extra legal renewals.}
\]

The evidence is strongest when the buyer mechanism is exact SIRET or direct
same-SIREN and weakest when the link depends on historical alias
propagation plus low text similarity or low top1--top2 margin.

\section{Procurement-Profile Domain Audit}
\label{sec:profile}

\subsection{Role in the pipeline}

The profile-domain audit is a diagnostic overlay after M1. It answers:
\emph{when M1 says that two buyer-name variants belong together, do
procurement-profile domains provide support, neutral shared-platform
evidence, weak support, ambiguity, or no usable evidence?}

It does not:

\begin{itemize}
  \item create a SIREN or SIRET;
  \item create a buyer key;
  \item create a candidate pair;
  \item increase the link count;
  \item replace M0 or M1.
\end{itemize}

\subsection{Domain classification}

URLs are normalized into domains and registered domains. Known shared
platforms such as \texttt{achatpublic.com},
\texttt{marches-publics.gouv.fr}, \texttt{marches-securises.fr},
\texttt{boamp.fr}, and related platforms are classified as generic. Other
domains are classified from concentration patterns:

\begin{itemize}
  \item \texttt{GENERIC\_PLATFORM}: broad multi-buyer platforms;
  \item \texttt{BUYER\_SPECIFIC}: one buyer, enough repeated evidence;
  \item \texttt{HIGHLY\_CONCENTRATED}: few buyers with one dominant buyer;
  \item \texttt{SHARED\_INSTITUTIONAL}: limited shared institutional use;
  \item \texttt{INSUFFICIENT\_DATA}, \texttt{AMBIGUOUS}, or
        \texttt{MALFORMED}.
\end{itemize}

The audit found 1,408 generic-platform domains and 433 buyer-specific or
highly concentrated domains.

\subsection{Evidence status}

\begin{table}[h]
\centering
\caption{Profile evidence over M1 aliases and incremental links.}
\begin{tabular}{lrr}
\toprule
Profile evidence status & Alias rows & Incremental links\\
\midrule
\texttt{NO\_PROFILE\_EVIDENCE} & 2,423 & 76\\
\texttt{NEUTRAL\_SHARED\_PLATFORM} & 593 & 59\\
\texttt{STRONG\_SUPPORT} & 110 & 86\\
\texttt{AMBIGUOUS} & 29 & 2\\
\texttt{WEAK\_SUPPORT} & 4 & 50\\
\bottomrule
\end{tabular}
\end{table}

For the 1,247 auto-propagation aliases specifically, the profile evidence
counts are: 928 no profile evidence, 253 neutral shared platform, 53 strong
support, 12 ambiguous, and 1 weak support. For the 198 historical-alias
incremental links, the counts are: 68 strong support, 54 neutral shared
platform, 41 weak support, 34 no profile evidence, and 1 ambiguous.

\subsection{Profile-audited variant}

The profile-audited variant only removes explicit profile conflicts among
historical-alias candidate pairs. It does not add new pairs. In the current
run it removes zero links:

\begin{table}[h]
\centering
\caption{M1 versus profile-audited M1.}
\begin{tabular}{lrr}
\toprule
Metric & M1 & Profile-audited M1\\
\midrule
Eligible sources & 3,159 & 3,159\\
Sources with candidates & 1,504 & 1,504\\
Candidate pairs & 8,228 & 8,228\\
Accepted links & 847 & 847\\
Blocking coverage & 47.6\% & 47.6\%\\
Linking rate & 26.8\% & 26.8\%\\
Candidate reuse rate & 29.5\% & 29.5\%\\
Jaccard versus M1 & 1.000 & 1.000\\
Links removed versus M1 & 0 & 0\\
\bottomrule
\end{tabular}
\end{table}

\section{Workflow Diagrams}
\label{sec:figures}

\begin{figure}[h]
\centering
\includegraphics[width=\linewidth]{m1_enrichment_profile_pipeline.png}
\caption{M1 enrichment and profile-domain audit workflow.}
\end{figure}

\begin{figure}[h]
\centering
\includegraphics[width=0.95\linewidth]{m1_reproducible_runflow.png}
\caption{Reproducible run flow for the M1 experiment and audit outputs.}
\end{figure}

\section{Consistency Checks}
\label{sec:checks}

The M1 consistency audit contains ten checks, all passing:

\begin{itemize}
  \item no accidental many-to-many notice merge;
  \item no duplicated notice IDs in M1 sources;
  \item no silent overwrite of BOAMP SIRET;
  \item every enriched identifier has provenance;
  \item all 1,002 conflict rows are traceable;
  \item every accepted M1 link exists in the candidate table;
  \item every event row has a valid candidate;
  \item censored rows do not contain accepted candidates;
  \item event and censoring times are positive;
  \item M0 reference files are present for hash trace.
\end{itemize}

The profile-audit consistency layer contains fourteen checks, also all
passing. The most important are that profile domains are not used as legal
identifiers, generic platforms do not create buyer matches, explicit
conflicts are traceable, every accepted audited link exists in the
candidate table, and M0/M1 outputs remain read-only reference inputs.

\section{Limitations}
\label{sec:limitations}

M1 improves coverage, but it introduces additional assumptions. The main
limitations are:

\begin{enumerate}
  \item \textbf{External-enrichment coverage starts in 2024.} Direct
        notice-level enrichment does not cover the full 2015--2026 study
        period, so historical propagation is needed for older notices.
  \item \textbf{Identifier conflict is real.} The audit finds 1,002 cases
        where BOAMP SIRET-derived SIREN and enriched SIREN disagree. Those
        cases are flagged, but they show why enrichment cannot be treated
        as automatically superior to BOAMP.
  \item \textbf{Historical aliases are plausible, not legal proof.} The
        automatic alias bridge requires unique SIREN evidence by normalized
        name and department, but it is still a heuristic.
  \item \textbf{Incremental links have weaker text evidence.} The median
        text similarity of incremental links is 0.078, much lower than the
        unchanged M0 links at 0.175.
  \item \textbf{Profile-domain evidence is incomplete and name-based.} It
        has no notice identifiers and no legal identifiers, so it can audit
        or flag support/conflict but cannot identify buyers.
  \item \textbf{The event is still a proxy recurrence.} M1 changes the
        constructed candidate set; it does not observe legal contract
        renewal status.
\end{enumerate}

\section{Recommended Interpretation}
\label{sec:interpretation}

The correct reading is:

\[
\begin{aligned}
&\text{M0 is conservative and official;}\\
&\text{M1 is wider and informative;}\\
&\text{manual validation decides whether M1 can become primary.}
\end{aligned}
\]

M1 demonstrates that buyer-name fragmentation materially suppresses M0
candidate coverage. However, the additional links should be reported as a
sensitivity result, not pooled with M0 as ground truth. The best immediate
use of M1 is to guide manual validation: prioritize incremental links from
historical aliases, low-margin links, CPV mismatches, and rows with profile
ambiguity or missing profile evidence.

\section{Reproducibility}
\label{sec:repro}

The M1 enrichment and audit outputs are regenerated by:

\begin{verbatim}
python3 scripts/build_m1_buyer_siren_experiment.py
python3 scripts/build_m1_profile_domain_audit.py
MPLCONFIGDIR=/tmp/matplotlib-m1-report \
  python3 scripts/make_m1_academic_diagrams.py
\end{verbatim}

The consolidated report itself is rebuilt from \texttt{reports/} with:

\begin{verbatim}
latexmk -cd -pdf -interaction=nonstopmode \
  -halt-on-error boamp_m1_technical_report.tex
\end{verbatim}

\section{Relation to Previous M1 Notes}
\label{sec:relation}

This document consolidates the previous M1 Markdown notes into one
canonical report. The short generated notes separated the SIREN enrichment
experiment, the procurement-profile audit, and the integrated diagram
summary. This report keeps those parts together so the reader can follow
the full chain:

\[
\begin{aligned}
\text{M0 reference}
&\rightarrow \text{SIREN enrichment}
\rightarrow \text{M1 candidate construction}\\
&\rightarrow \text{M0/M1 comparison}
\rightarrow \text{profile-domain audit}\\
&\rightarrow \text{interpretation and limits.}
\end{aligned}
\]

\section{Conclusion}
\label{sec:conclusion}

M1 is technically useful and scientifically important because it shows that
buyer-identity repair changes the observable recurrence signal. On the
same eligible source population and the same threshold, it raises blocking
coverage from 39.1\% to 47.6\% and accepted proxy-recurrence links from
618 to 847. That is a material sensitivity result.

The conservative conclusion is equally important: M1 should not replace M0
yet. The largest incremental channel is historical alias reconciliation,
the incremental links have lower median text similarity than unchanged M0
links, and profile-domain data are audit evidence rather than legal
identifiers. The current repository position is therefore methodologically
consistent: report M0 as the official BOAMP-only baseline, report M1 as a
buyer-enriched sensitivity analysis, and use the exported manual-validation
samples to decide whether any future report can promote an enriched
specification.

\end{document}
